% !PRINTED 484
\weberpage{484}

Geht nämlich durch $\pi$ der Coëfficient $A_{h,i,k}$ in $A_{h,i,k}'$ über, so ist vermittelst der Transformation $(y) = A(x)$:
\[
\sum A_{h,i,k}' \, y_h y_i y_k = \sum A_{h,i,k} \, x_h x_i x_k,
\]
und vermittelst der Substitution $\xi = A_1(\eta)$ besteht die Identität $f(\xi_1, \xi_2, \xi_3) = f(\eta_1, \eta_2, \eta_3)$. Folglich ergiebt sich nach dem Satze \S.\,129, 3.:
\[
\sum A_{h,i,k}' \, \frac{\partial^3 f}{\partial\eta_h\partial\eta_i\partial\eta_k} = \sum A_{h,i,k} \, \frac{\partial^3 f}{\partial\xi_h\partial\xi_i\partial\xi_k}.
\]

Bedeutet also $\pi$ irgend eine Permutation der Gruppe $P_{168}$, durch die $X_1, X_2, X_3$ in $Y_1, Y_2, Y_3$ übergeht, und
\[
A = \begin{pmatrix} a_1, & b_1, & c_1\\ a_2, & b_2, & c_2\\ a_3, & b_3, & c_3 \end{pmatrix}
\]
die entsprechende Substitution aus der Gruppe $G_{168}$, so haben wir zu setzen:
\[
\begin{aligned}
\xi_1 &= a_1\eta_1 + a_2\eta_2 + a_3\eta_3\\
\xi_2 &= b_1\eta_1 + b_2\eta_2 + b_3\eta_3\\
\xi_3 &= c_1\eta_1 + c_2\eta_2 + c_3\eta_3,
\end{aligned}
\]
und die Invarianteneigenschaft von $L$ giebt die Relation
\weq{10}{Y_1\eta_1 + Y_2\eta_2 + Y_3\eta_3 = X_1\xi_1 + X_2\xi_2 + X_3\xi_3,}
oder entwickelt
\weq{11}{\begin{aligned} Y_1 &= a_1X_1 + b_1X_2 + c_1X_3\\ Y_2 &= a_2X_1 + b_2X_2 + c_2X_3\\ Y_3 &= a_3X_1 + b_3X_2 + c_3X_3, \end{aligned}}
d.\,h. die Permutation $\pi$, auf die Functionen $X_1, X_2, X_3$ angewandt, hat denselben Erfolg, wie die lineare Substitution $A(X_1, X_2, X_3)$, und demnach sind die $X_1, X_2, X_3$ solche Functionen, wie sie unser Problem verlangt (Schluss des \S.\,128).

\begin{center}
\S. 131.

\gesperrt{Möglichkeit der Bestimmung der Functionen $X_1, X_2, X_3$.}
\end{center}

Um die Zurückführung der Gleichung 7\textsuperscript{ten} Grades mit der Gruppe $P_{168}$ auf das Formenproblem der Gruppe $G_{168}$ vollständig sicher zu stellen, bleibt noch Eines übrig:
%CONT>>
